\documentclass[12pt, letterpaper]{article}
\date{\today}
\usepackage[margin=1in]{geometry}
\usepackage{amsmath}
\usepackage{hyperref}
\usepackage{cancel}
\usepackage{amssymb}
\usepackage{fancyhdr}
\usepackage{pgfplots}
\usepackage{booktabs}
\usepackage{pifont}
\usepackage{amsthm,latexsym,amsfonts,graphicx,epsfig,comment}
\pgfplotsset{compat=1.16}
\usepackage{xcolor}
\usepackage{tikz}
\usetikzlibrary{shapes.geometric}
\usetikzlibrary{arrows.meta,arrows}
\newcommand{\Z}{\mathbb{Z}}
\newcommand{\N}{\mathbb{N}}
\newcommand{\R}{\mathbb{R}}
\newcommand{\Q}{\mathbb{Q}}
\newcommand{\C}{\mathbb{C}}
\newcommand{\F}{\mathbb{F}}

\newcommand{\Po}{\mathcal{P}}
\newcommand{\Pro}{\mathbb{P}}
\author{Alex Valentino}
\title{501 homework}
\pagestyle{fancy}
\renewcommand{\headrulewidth}{0pt}
\renewcommand{\footrulewidth}{0pt}
\fancyhf{}
\rhead{
	Homework 7\\
	501
}
\lhead{
	Alex Valentino\\
}
\begin{document}
\begin{enumerate}
	\item[7] Let $\{f_n\}$ be a sequence in $L^1(X, \mathcal{M},\mu)$ such that $f_n \to f$ uniformly, $\mu(X) < \infty$.
	We want to show that $f \in L^1$ and $\lim \|f - f_n\|_1 = 0$.  Note that for some $\epsilon > 0$ there exists 
	$N_\epsilon$ such that for all $n \geq N_\epsilon$ we have that $|f - f_n| < \frac{\epsilon}{\mu(X)}$.  
	Therefore $\int_X |f - f_n| d\mu < \frac{\epsilon}{\mu(X)} \int_x d \mu = \epsilon$.  Therefore in the limit it goes 
	to zero.  Additionally $\int_X |f| d \mu \leq \int_X |f - f_n| d \mu + \int_X |f_n| d \mu < \infty + \epsilon =
	 \infty$, giving us that $f \in L^1$.  Note that this fails to hold if $\mu(X) = \infty$.  In particular if one 
	 takes the functions $f_n(x) = \sum_{k=1}^\infty \frac{1}{k^{1+\frac{1}{n}}}1_{(k-1,k)}(x)$, each $\int_\R |f_n| d \mu$ is guarenteed to converge by the p-series tests.  However if one takes $n$ to infinity you get the harmonic series,
	 which diverges.  
	
	\item[8] Let $f(x) = x^{-1/2}1_{(0,1)}, \{q_n\}_{n \in \N}$ be an enumeration of the rationals, and 
	$g(x) = \sum_{n=1}^\infty f(x - q_n)$.  We claim that $g \in L^1$.  Note that first each $2^{-n}f(x - q_n)$ is an 
	integrable function, therefore we can interchange the  sum and the integral:
	$\int_\R |g| d \mu = \int_\R \sum_{n=1}^\infty 2^{-n} f(x-q_n) d\mu = \sum_{n=1}^\infty \int_\R 2^{-n}f(x-q_n) d\mu$.
	Furthermore, by the construction of $f,$ $\int_\R f(x-q_n) d \mu = \int_\R f d \mu$, since the shifting of both the function and the indidcator exactly cancel out.  Therefore $\int_\R |g| = \sum_{n=1}^\infty 2^{-n}\int_\R f d \mu = 
	\sum_{n=1}^\infty 2^{1-n} = 1$.  We additionally claim for all intervals $I \subset \R, \lambda \geq 1$ that
	$\mu(I \cap \{x: g(x) > \lambda\}) > 0$.   Let $(a,b) = I, \lambda \geq 1$, and let $q_a, q_b$ be rational numbers such that 
	$a < q_a < q_b < b$  and $q_b - q_a < 1$.  If we then consider $2^{-\alpha}f(x - \frac{q_b + q_a}{2})$ where 
	$\alpha$ is the index of $\frac{q_b + q_a}{2}$ we are guaranteed that $g$ has a singularity contained within said 
	interval.  Furthermore, the interval $(\frac{q_b + q_a}{2}, \min(\frac{1}{\frac{q_b + q_a}{2} +\lambda^2}, q_b)$ is
	guaranteed to be of non-zero measure.  Thus $\mu(I \cap \{x: g(x) > \lambda\}) > 0$.  Note that since $g$ is finite
	almost everywhere, then on that same set $g^2$ is finite.  However by cauchy-schwarz we have that 
	$$
     \sum_{n=1}^\infty 4^{-n}f(x-q_n)^2  \leq g^2.
	$$
	Since for a single $n$, we have that $\int_\R f(x - q_n)^2 = \int_0^1 \frac{1}{x} dx = \infty$, and since the rationals are dense in $\R$ then every interval will necessarily contain a term which diverges, thus making it not integrable on said interval.
	\item[9] Let $f_n(x) = 1_{E_n}(x)$ be a cauchy sequence of indicators in $L^1(X,\mathcal M, \mu)$.  Note by the completion of 
	$L^1$ then there exists $f \in L^1$ such that $\lim f_n = f$ almost everywhere.  Note that since the set $E$ where they 
	disagree is a set of measure zero, we can simply multiple $f$ by $1_{E^c}$, which ensures that $f 1_{E^c}$  is an 
	indicator function.  As on $E^c, f$ is an indicator function since it can only be $1$ or $0$.  On $E$ if we set 
	all values to $0$ then it still is an indicator function.  Since they differ only on a set of measure 0 then the 
	sequence converges to an indicator function in $L^1$.
	\item[10]
	\item[11] Let $t_n$ be some sequence converging to $t > 0$. Note that $f(x)g(t_n x) \leq M|f(x)|$ for all $x \in \R$ and $\int_\R M |f(x)|d \mu < \infty$.  Thus by LDCT we have that $\lim F(t_n) = \int_\R \lim f(x) g(t_n x) d\mu$.  
	Furthermore any constant postitive "multiple" of the set is still of measure zero.  Therefore $\lim g(t_nx) = g(tx)$.
	Thus $\lim F(t_n) = F(t)$.  Thus $F$ is continuous.
	 
	\item[12] Let $(X,\mathcal M, \mu)$ be a finite measure space, $\{f_n\}$ be a sequence of integrable functions such that 
	$f_n \to 0$ in measure and $\sup \| |f_n|^2 \|_1 < \infty$.  We claim that the sequence converges to 0 in $L^1$.
	Let $\epsilon > 0$ and $E_{n,\epsilon} = \{x: |f(x)| \geq \epsilon\}$.  Note that $\lim \mu(E_{n,\epsilon}) = 0$ 
	for all $\epsilon$ and that  
	\begin{align*}
	|\int_X f_n d \mu|^2 &\leq |\int_{E_{n,\epsilon}} f_n|^2 + |\int_{E_{n,\epsilon}^c} f_n|^2 +|\int_{E_{n,\epsilon}} f_n||\int_{E_{n,\epsilon}^c} f_n|\\
	&\leq |\int_{E_{n,\epsilon}} f_n|^2 + \epsilon \mu(X)|\int_{E_{n,\epsilon}} f_n| + \epsilon^2\mu(X) \\
	&\leq \mu(E_{n,\epsilon})\int_X |f_n|^2 d \mu + \epsilon \mu(X)|\int_{E_{n,\epsilon}} f_n|+ \epsilon^2\mu(X)^2
	\end{align*}
	Since $\sup \||f_n|^2\|_1 < \infty$, then let $M = \sup \||f_n|^2\|_1$.  Additionally, since $\lim 
	\mu(E_{n,\epsilon}) = 0$ then $\lim \int_{E_{n,\epsilon}} f_n d\mu = 0$ as an integral on a set of measure zero is 
	zero.  Therefore the inequality is bounded by $\epsilon^2 \mu(X)^2$, therefore the integral goes to zero.
	
	\item[13] Let $(X, \mathcal{M}, \mu)$ be a measure space, $\mu$ is a semifinite measure, $f$ is a non-negative 
	function measurable function such that for all $g \in L^+(X,\mathcal{M}, \mu)$, 
	$$
	\int_X fg d\mu < \infty.
	$$
	We claim that $f$ is uniformly bounded almost everywhere.  Suppose not.  Then for all $M > 0$ there exists 
	a set $E_M$ such that for all $x \in E_M, f(x) > M$
\end{enumerate}
\end{document}
